
% ============================================================
%  INTRODUCTION GÉNÉRALE
% (pagenumbering/setcounter are in main.tex)
% ============================================================
\cleardoublepage

\chapter*{Introduction Générale}
\addcontentsline{toc}{chapter}{Introduction Générale}
\markboth{Introduction Générale}{Introduction Générale}

\section*{Contexte mondial}
L'agriculture mondiale traverse une mutation profonde sous l'impulsion de
la \textit{quatrième révolution industrielle}. Les technologies de l'Internet
des Objets (IoT), de l'intelligence artificielle (IA) et des systèmes
d'information géographique (GIS) convergent pour donner naissance au
concept de \textbf{Smart Farming} — une agriculture de précision, connectée
et pilotée par la donnée. Selon la FAO~\cite{fao2023}, la population mondiale
devrait atteindre 9,7 milliards d'individus en 2050, imposant une augmentation
de 70\% de la production alimentaire mondiale. Dans ce contexte, l'IA
agricole représente un levier stratégique incontournable.

\section*{Enjeux tunisiens}
La Tunisie, dont le secteur agricole représente environ 10\% du PIB et
emploie 16\% de la population active (UTAP~\cite{utap2023}), fait face
à des défis structurels majeurs : morcellement foncier, vieillissement de la
population agricole, accès limité aux technologies et dépendance vis-à-vis
de solutions étrangères coûteuses inadaptées au contexte local. L'absence
d'outils numériques en langue locale (Darija tunisienne) constitue un
obstacle particulièrement saillant pour l'adoption technologique par les
éleveurs.

\section*{Présentation du projet}
\textbf{Smart Farm AI v3.0 Enterprise} est une plateforme d'agriculture
intelligente conçue et développée en réponse à ces défis. Elle couvre la
gestion multi-espèces (bovins, ovins, caprins, volailles, abeilles), l'analyse
d'image par vision par ordinateur, un assistant IA conversationnel en Darija
tunisienne, la surveillance IoT en temps réel et un ERP avicole complet.
L'ensemble est déployé selon une architecture multi-tenant sécurisée, accessible
sur mobile via une PWA offline-first.

\section*{Objectifs du mémoire}
Ce mémoire poursuit quatre objectifs complémentaires :
\begin{enumerate}
  \item Documenter la démarche de conception et d'implémentation d'une
        architecture logicielle complexe intégrant IA, IoT et web ;
  \item Évaluer les performances des modèles ML déployés selon des méthodes
        statistiques rigoureuses (intervalles de confiance, validation croisée) ;
  \item Analyser la validité interne et externe des résultats obtenus ;
  \item Proposer une contribution à l'état de l'art sur les systèmes RAG+LLM
        en langue dialectale non-standard (Darija).
\end{enumerate}

\section*{Hypothèses de recherche}
Ce travail pose trois hypothèses testables, vérifiées expérimentalement dans les
chapitres 3 et 4 :
\begin{description}
  \item[\textbf{H1} — Efficacité Darija] Un pipeline RAG+LLM local (Labess-7B) obtiendra
        un accord inter-annotateurs supérieur à $\kappa = 0.60$ (accord substantiel,
        Landis \& Koch~\cite{landis1977}) pour des requêtes en Darija tunisienne
        dans le domaine de l'élevage. \textit{Vérifiée} : $\kappa = 0.76$ (Chap.~3).
  \item[\textbf{H2} — Stabilité YOLO] Trois runs indépendants de YOLOv8 avec des
        graines aléatoires différentes produiront un mAP@50 global dont l'écart-type
        est inférieur à $\pm 1\%$. \textit{Vérifiée} : $\sigma = 0.45\%$ (Chap.~3).
  \item[\textbf{H3} — Adoption PWA] Une application mobile offline-first atteindra
        un taux d'adoption supérieur à $80\%$ parmi les travailleurs agricoles formés
        après une semaine d'usage sur les 7 fermes pilotes.
        \textit{Vérifiée} : 91\% d'adoption (Chap.~4).
\end{description}

\section*{Contributions originales}
Par rapport aux travaux existants, ce projet apporte trois contributions nouvelles :
\begin{enumerate}
  \item \textbf{Première évaluation quantitative RAG+LLM en Darija agricole} :
        à notre connaissance, aucun travail publié ne combine RAG + LLM + évaluation
        BLEU/ROUGE/Kappa pour le dialecte tunisien dans un contexte d'élevage.
  \item \textbf{Protocole de validation statistique pour un ERP industriel} :
        l'utilisation de la validation croisée 5-fold avec \textit{TimeSeriesSplit},
        d'intervalles de confiance à 95\% et d'une analyse des menaces à la validité
        est inhabituellement rigoureuse pour un système de gestion agricole opérationnel.
  \item \textbf{Architecture open-source agro-tech souveraine} : l'intégration de
        11 catégories YOLO, 5 modèles ML avicoles, RAG ChromaDB et LLM local dans
        un déploiement Docker entièrement open-source constitue un patron d'architecture
        reproductible pour l'agri-tech en pays émergents à connectivité limitée.
\end{enumerate}

\section*{Structure du document}
Le mémoire est organisé en quatre chapitres. Le \textbf{Chapitre~1} présente
l'organisme d'accueil, le contexte du projet et la méthodologie de gestion
de projet hybride CRISP-DM/Scrum. Le \textbf{Chapitre~2} détaille l'analyse
des besoins et l'architecture globale du système. Le \textbf{Chapitre~3}
expose la modélisation des composants IA et l'évaluation quantitative des
performances. Le \textbf{Chapitre~4} décrit la réalisation technique et
la validation industrielle sur le terrain.

% ============================================================
%  CHAPITRE 1 : CONTEXTE GÉNÉRAL ET CADRE MÉTHODOLOGIQUE
% ============================================================
